We list a few open questions that follow directly from our work.

\begin{itemize}
    \item In \cite{Eppstein1990}, it was asked if a word of rank one for a total DFA can be found in \NC. Similarly, can a matrix of minimum rank for a total DFA be computed in \NC?

    \item Given an unambiguous morphism $M \colon \Sigma \to \{0, 1\}^{Q \times Q}$ and a vector $\alpha \in \mathbb{Q}_{>0}^Q$, can a basis of  $\spn{\alpha^T M(w) \mid w \in \Sigma^*}$ be computed faster than in $\OO(|Q|^3)$ time? This would improve algorithms for several fundamental problems for weighted automata~\cite{KieferWeighted20}. Similarly, for total DFAs, computing a basis of $U$ from \autoref{subsec:pseudocolumns} in subcubic time (see the proof of \autoref{lem:pseudcolumn-time}) would allow to compute the rank of a total DFA faster than in cubic time.

    \item Can one decide if a total DFA has rank one in strongly subquadratic time (in the number of states)? This seems to be a major open problem in the area of synchronising automata.
    
    \item The bottleneck in the time complexity in \autoref{thm:linalg-time} is the very first step, computing $\Mer(q)$ via digraph search in the square digraph of $\Aa$. The number of edges of this digraph can be quadratic in the number of its vertices~\cite[Appendix A]{Kiefer2019Arxiv}, hence of order~$|Q|^4$. Can $\Mer(q)$ be computed faster than in time $\OO(|Q|^4)$? Very little seems to be known about general properties of square automata of DFAs or UFAs.

    \item Finally, a natural continuation of this work is to consider the minimum nonzero rank of zero-one matrix monoids. It is equal to the rank of each nonzero matrix in the 0-minimal ideal of the monoid. It is not known how to compute it in \NC even for total DFAs. The main motivation once again comes from the degree of variable-length codes, see \cite[Chapter 9]{Berstel2010} for more details. 
\end{itemize}